\section{The aspect ratio formula}\label{sec:family}

An \emph{iterated function system} (IFS) is a finite collection
of contractions $f_1,\ldots,f_N$ on a complete metric space.
By Hutchinson's theorem \cite{Hutchinson1981} there exists a unique
non-empty compact set, the \emph{attractor}, satisfying
$\mathcal{A} = f_1(\mathcal{A}) \cup \cdots \cup f_N(\mathcal{A})$.

Consider an IFS
\[
  f_k(\mathbf{v}) = A R_k(\mathbf{v} + \mathbf{t}_k), \qquad k=1,\ldots,N,
\]
where $A$ is a $2\times 2$ real matrix with complex eigenvalues
$\lambda,\bar\lambda\in\mathbb{C}\setminus\mathbb{R}$, $|\lambda|<1$,
and each $R_k$ is a real $2\times 2$ matrix that is an isometry
\emph{in the Jordan basis of $A$}: that is, $P^{-1}R_kP$ acts as
multiplication by some $\varepsilon_k\in\mathbb{C}$, $|\varepsilon_k|=1$,
where $P$ is the real matrix whose columns are $\mathrm{Re}(\mathbf{p})$
and $\mathrm{Im}(\mathbf{p})$ for an eigenvector $\mathbf{p}$ of $A$
for eigenvalue $\lambda$.
In the Jordan basis, writing $t_k' = (P^{-1}\mathbf{t}_k)^{(1)}
+ i(P^{-1}\mathbf{t}_k)^{(2)} \in \mathbb{C}$,
the centering condition is $\sum_{k=1}^N \varepsilon_k t_k' = 0$.

Define the \emph{digit variance} $\sigma^2 := \tfrac{1}{N}\sum_k|t_k'|^2 > 0$,
the \emph{weighted algebraic variance}
$\tau_\varepsilon := \tfrac{1}{N}\sum_k \varepsilon_k^2(t_k')^2 \in \mathbb{C}$,
the \emph{weight second moment} $\alpha_2 := \tfrac{1}{N}\sum_k \varepsilon_k^2 \in \mathbb{C}$,
and the \emph{weighted digit anisotropy}
\[
  \kappa_\varepsilon := \frac{|\tau_\varepsilon|}{\sigma^2} \in [0,1].
\]

\begin{lemma}\label{lem:kappa-invariant}
Under the Jordan basis change $\mathbf{t}_k \mapsto t_k' = P^{-1}\mathbf{t}_k$
(identified as a complex number), the digit anisotropy is preserved:
if we define $\kappa$ using the original digits $\mathbf{t}_k$ identified
as complex numbers $t_k$ via the standard basis, and $\kappa_\varepsilon$ using
the Jordan-basis digits $t_k'$, then when all $\varepsilon_k = 1$,
\[
  \kappa'
  := \frac{|\sum_k (t_k')^2|}{\sum_k |t_k'|^2}
   = \kappa
  := \frac{|\sum_k t_k^2|}{\sum_k |t_k|^2}.
\]
\end{lemma}

\begin{proof}
The columns of $P$ are $\mathrm{Re}(\mathbf{p})$ and $\mathrm{Im}(\mathbf{p})$,
so $P^{-1}$ acts on $\mathbb{R}^2\cong\mathbb{C}$ as multiplication by
$1/p$, where $p\in\mathbb{C}$ is the eigenvector regarded as a complex number.
Therefore $t_k' = t_k/p$, and
\[
  \kappa'
  = \frac{|\sum_k (t_k/p)^2|}{\sum_k |t_k/p|^2}
  = \frac{|\sum_k t_k^2|/|p|^2}{\sum_k |t_k|^2/|p|^2}
  = \kappa. \qedhere
\]
\end{proof}

\begin{theorem}\label{thm:family}
Let $A$ be a $2\times 2$ real matrix with eigenvalues
$\lambda,\bar\lambda\in\mathbb{C}\setminus\mathbb{R}$, $|\lambda|<1$,
and let $\{f_k(\mathbf{v})=AR_k(\mathbf{v}+\mathbf{t}_k)\}_{k=1}^N$
be an equal-weight IFS where each $R_k$ is an isometry in the Jordan
basis of $A$ with corresponding complex factor $\varepsilon_k$,
$|\varepsilon_k|=1$.
Suppose the centering condition $\sum_k \varepsilon_k t_k' = 0$ holds
in the Jordan basis and $\sigma^2>0$.
Then the invariant measure has aspect ratio
\begin{equation}\label{eq:family}
  \mathrm{AR}^2 = \frac{|1-\lambda^2\alpha_2| - (1-|\lambda|^2)\,\kappa_\varepsilon}
                      {|1-\lambda^2\alpha_2| + (1-|\lambda|^2)\,\kappa_\varepsilon},
\end{equation}
where $\alpha_2 = \tfrac{1}{N}\sum_k\varepsilon_k^2$,
$\tau_\varepsilon = \tfrac{1}{N}\sum_k\varepsilon_k^2(t_k')^2$,
and $\kappa_\varepsilon = |\tau_\varepsilon|/\sigma^2$.
\end{theorem}

\begin{proof}
By Lemma~\ref{lem:kappa-invariant} it suffices to work in the Jordan
basis of $A$, where $f_k(z) = \lambda\varepsilon_k(z + t_k')$
with $t_k'\in\mathbb{C}$.
Let $Z$ be distributed according to the invariant measure $\mu$.
The fixed-point equation
\begin{equation}\label{eq:fp}
  E[h(Z)] = \frac{1}{N}\sum_{k=1}^N E\bigl[h\bigl(\lambda\varepsilon_k(Z + t_k')\bigr)\bigr]
\end{equation}
applied to $h(z)=z$ gives, using $\sum_k\varepsilon_k t_k'=0$,
\[
  E[Z] = \lambda\alpha_1 E[Z],
  \qquad \alpha_1 := \tfrac{1}{N}\textstyle\sum_k\varepsilon_k,
\]
so $E[Z]=0$ (since $|\lambda\alpha_1|\leq|\lambda|<1$).
Applying~\eqref{eq:fp} to $h(z)=z^2$:
\[
  E[Z^2] = \lambda^2\alpha_2\,E[Z^2] + \lambda^2\tau_\varepsilon,
  \qquad\text{so}\qquad
  E[Z^2] = \frac{\lambda^2\tau_\varepsilon}{1-\lambda^2\alpha_2}.
\]
Applying~\eqref{eq:fp} to $h(z)=|z|^2$, noting $|\varepsilon_k|=1$:
\[
  E[|Z|^2] = |\lambda|^2 E[|Z|^2] + |\lambda|^2\sigma^2,
  \qquad\text{so}\qquad
  E[|Z|^2] = \frac{|\lambda|^2\sigma^2}{1-|\lambda|^2}.
\]
The ellipticity ratio is
\[
  \rho = \frac{|E[Z^2]|}{E[|Z|^2]}
       = \frac{(1-|\lambda|^2)\,\kappa_\varepsilon}{|1-\lambda^2\alpha_2|}.
\]
The standard identity $\mathrm{AR}^2=(1-\rho)/(1+\rho)$
yields~\eqref{eq:family}.
\end{proof}

\begin{corollary}\label{cor:scalar}
When $A = aI$ for $a\in\mathbb{C}$, $|a|<1$, and all $R_k = I$
(so $\varepsilon_k=1$ and $\alpha_2=1$), the Jordan basis coincides
with the standard basis, $\lambda=a$, $\kappa_\varepsilon = \kappa$,
and~\eqref{eq:family} reduces to
\[
  \mathrm{AR}^2 = \frac{|1-a^2| - (1-|a|^2)\,\kappa}
                      {|1-a^2| + (1-|a|^2)\,\kappa},
\]
which is the aspect ratio formula for the scalar IFS
$\{f_k(z)=a(z+t_k)\}$ over $\mathbb{C}$.
\end{corollary}

\begin{remark}\label{rem:special-cases}
Theorem~\ref{thm:family} unifies several previously separate results.
The following table summarises the main special cases.

\medskip
\begin{center}
\begin{tabular}{cccp{4.5cm}}
\hline
$A$ & $R_k$ & $\varepsilon_k$ & Result \\
\hline
$aI$ & $I$ & $1$ & scalar formula (Corollary~\ref{cor:scalar}) \\
$aI$ & $\pm I$ & $\pm 1$ & sign-flip (Remark~\ref{rem:sign-flip}) \\
$aI$ & rotations & $e^{i\theta_k}$ & rotation-weights case \\
general $A$ & Jordan-basis isometries & $|\varepsilon_k|=1$ & Theorem~\ref{thm:family} \\
\hline
\end{tabular}
\end{center}
\medskip

Three further special cases of the general formula are noteworthy:
\begin{enumerate}[(a)]
\item \textbf{Sign flips} ($\varepsilon_k=\pm 1$): then $\alpha_2=1$ and
  $\kappa_\varepsilon = |\sum_k \varepsilon_k^2(t_k')^2|/\sigma^2$,
  recovering Remark~\ref{rem:sign-flip}.
\item \textbf{Uniform rotation} ($\varepsilon_k = e^{i\theta}$ for all $k$):
  then $\alpha_2 = e^{2i\theta}$ and $\tau_\varepsilon = e^{2i\theta}\tau$,
  so $\kappa_\varepsilon = \kappa$ and the AR is unchanged — the attractor
  is merely rotated by $\theta$.
\item \textbf{$N$-th roots of unity} ($\varepsilon_k = e^{2\pi i k/N}$,
  $N\geq 3$): then $\alpha_2 = 0$
  and~\eqref{eq:family} reduces to
  $\mathrm{AR}^2 = (1-(1-|\lambda|^2)\kappa_\varepsilon)/(1+(1-|\lambda|^2)\kappa_\varepsilon)$.
\end{enumerate}
\end{remark}